%% parts/p07/arithmetic-proof-theory-supplement.tex
% Necessity: the exact inherited ordinal-analysis section precedes but does
% not contain the separately attributed D-P8 calibration tables and new
% theory-presentation, Veblen, and finitism apparatus.
% Definition: this separately authored derivative retains the historical
% exposition corresponding to lines 7--46 and 65--95, replaces the former
% N-a-049 stub at line 110 by a selected propositional ordinal assignment and
% annex proof, and revises the former lines 112--309 calibration material.
% Standing: it is an editorial reader supplement, not inherited source or an
% admitted local proof of Gentzen's full arithmetic consistency theorem.
% Nonredundancy: the duplicate Gentzen theorem and proof boundary at historical
% lines 48--63 and 97--108 are excluded; the exact source section controls them.
% Use: the Part VII driver retains the separate arithmetic chapter; this file
% is not represented as a verbatim restored source interval.

% BEGIN HISTORICAL SLICE 7-46
\chapter{Arithmetic Proof Theory}
\label{sec:p07-arithmetic-proof-theory}
\chaptermark{Arithmetic Proof Theory}

This chapter records what an ordinal analysis must contain before an ordinal
can be said to calibrate a formal theory.  The inherited source supplied an
honest cited boundary for Gentzen's analysis of arithmetic; it did not supply
the derivation coding, full assignment, or decrease proof.  That status is
preserved here.  The surrounding lookup material uses the ordinal arithmetic
and Cantor normal form developed in
Chapter~\ref{sec:p06-ordinal-arithmetic}.  Primary and modern sources include
\cite{Gentzen1936,Gentzen1938,TroelstraSchwichtenberg2000,
Takeuti1987,Schwichtenberg1977}.

\section{Gentzen's consistency analysis}
\label{sec:p07-gentzen-analysis}

\begin{hbmethod}{def:p07-gentzen-reduction-method}{Gentzen's reduction method}
For the selected sequent presentation of first-order arithmetic, a Gentzen
reduction analysis requires six separately specified objects:
\begin{enumerate}
  \item the exact arithmetic calculus, including its induction inferences;
  \item primitive-recursive codes for its finite derivations;
  \item a primitive-recursive notation system below \(\varepsilon_0\), with
        comparison and fundamental sequences;
  \item a primitive-recursive reduction operation on derivation codes;
  \item an ordinal-notation assignment and a case-by-case theorem that every
        critical reduction strictly lowers it; and
  \item the exact transfinite-induction schema used to exclude an infinite
        reduction, followed by extraction of the absence of a derivation of
        the empty sequent.
\end{enumerate}
The list is necessary because none of its members is determined by the bare
phrase ``ordinal assignment.''  It describes metalinguistic data attached to
one calculus, not additional axioms of arithmetic.  The reduction operation
uses the assignment in the cited termination theorem below; the notation
system alone proves no consistency statement.
\index{Gentzen reduction}
\index{ordinal assignment!derivations}
\end{hbmethod}
% END HISTORICAL SLICE 7-46

% BEGIN HISTORICAL SLICE 65-95
\begingroup\small
\begin{longtable}{@{}p{.20\textwidth}p{.36\textwidth}p{.35\textwidth}@{}}
\caption{Components of the cited Gentzen analysis.}
\label{tab:p07-gentzen-components}\\
\toprule
Component & Exact standing & Evidence retained in this handbook \\
\midrule
\endfirsthead
\toprule
Component & Exact standing & Evidence retained in this handbook \\
\midrule
\endhead
Arithmetic calculus & One named sequent calculus with equality and induction,
not an arbitrary presentation of \(\PA\) & Imported with the cited theorem;
the complete rule surface is not reproduced here \\
Derivation codes & Natural-number codes for finite proof trees, with
primitive-recursive local correctness & Imported; no local coding proof \\
Ordinal notations & Finite codes below \(\varepsilon_0\), with a
primitive-recursive comparison and fundamental-sequence operation & Imported;
Part~6 supplies ordinal meaning, not this coding algorithm \\
Reduction & Primitive-recursive transformation of the selected derivation
codes & Imported; no generic ``proof simplification'' is substituted \\
Assignment and decrease & A map from derivation codes to notation codes and a
strict decrease proof for every critical case & Imported; this is the
load-bearing theorem, not a diagrammatic arrow \\
Induction and extraction & The stated transfinite-induction schema excludes
infinite descent, after which normal-form analysis excludes the empty sequent
& Imported and relative to that metatheory \\
\bottomrule
\end{longtable}
\endgroup
% END HISTORICAL SLICE 65-95

% BEGIN HISTORICAL SLICE 110-110
\begin{hbdefinition}{def:p07-prop-cut-measure}{Ordinal assignment for the finite-set support presentation}
The \emph{finite-set support presentation} used here is defined from the
propositional restriction of
\hyperref[sheet:lk-defneg]{\texttt{sheet:lk-defneg}} as follows.  Antecedents
and succedents are finite sets.  Each propositional rule is the literal support
image of the corresponding displayed multiset rule.  Identity is contexted,
\(\Gamma\cup\{A\}\Rightarrow\Delta\cup\{A\}\); \(\bot L\) is contexted;
left and right weakening are combined into monotonicity under finite context
inclusion; exchange and contraction are identities of sets; Mix is absent.
Cut has the support form
\[
 \Gamma\Rightarrow\Delta\cup\{A\},\qquad
 \Pi\cup\{A\}\Rightarrow\Sigma
 \quad\leadsto\quad
 \Gamma\cup\Pi\Rightarrow\Delta\cup\Sigma,
\]
and is called \emph{clean} when
\(A\notin\Gamma\cup\Delta\cup\Pi\cup\Sigma\).  A nonclean Cut is erased:
one of its premises already weakens to its conclusion.  Thus only clean Cuts
enter the reduction below.

This presentation is necessary because the explicit multiset contraction rule
duplicates a Cut premise and invalidates the naive height-sum decrease.  It is
a formal presentation of the same propositional derivability relation, not a
new logic.  It differs nonredundantly from the registered sheet only by taking
supports and building exchange/contraction into the context representation.
The two translations used after the proof are exact: taking supports sends a
multiset derivation to this presentation, with contraction and exchange mapped
to identities and every multiset Mix mapped to support Cut on its Mix formula
(a nonclean image is immediately erased as above); choosing one occurrence of each set member sends a set
derivation back, using only weakening and contraction to reconcile unions.
The return translation introduces neither Cut nor Mix.  In the forward
direction, a source derivation containing Mix acquires support Cut as just
stated; a source derivation containing neither Cut nor Mix remains cut-free.

Put
\[
 r(p)=r(\bot)=0,
 \qquad
 r(B\mathbin\circ C)=1+\max\{r(B),r(C)\}
 \quad(\circ\in\{\land,\lor,\to\}).
\]
If any Cut node \(c=\operatorname{Cut}_A(D_0,D_1)\) has premise heights
\(h(D_0)\) and \(h(D_1)\), put
\(h(c)=h(D_0)+h(D_1)\).  For a derivation \(D\), define
\[
 o(D)=\mathop{\#}_{c\in\operatorname{Cut}(D)}
       \omega^{\omega r(A_c)+h(c)}<\omega^{\omega^2},
\]
where \(\#\) is Hessenberg natural sum and \(A_c\) is the Cut formula at
\(c\).

The reduction strategy is fixed, rather than left arbitrary.  If a nonclean
Cut remains, erase the first one in a fixed traversal order.  Otherwise let
\(r\) be the greatest rank of a Cut formula and choose a rank-\(r\) Cut whose
two premise subtrees contain no rank-\(r\) Cut; among such nodes the same
traversal chooses the first.  Thus duplicating a premise while commuting this
Cut over a branching inference can duplicate only lower-rank Cuts.  Explicit
contraction is absent, by definition of the presentation.
The assignment is a metalinguistic measure on this one finite derivation
grammar, not the \(\varepsilon_0\)-assignment for arithmetic.  It is used by
the complete decrease proof in Annex~\ref{annex:p7.arithmetic.1}.
\index{ordinal assignment!propositional cut reduction}
\end{hbdefinition}

\begin{center}
\begin{tabularx}{\textwidth}{@{}l X X@{}}
\toprule
Selected-Cut configuration & Reduction & Strict decrease \\
\midrule
one premise is the Cut formula's contexted identity & erase the Cut and retain the other premise by monotonicity & removes one summand \\
Cut formula not principal in one last inference & commute the selected Cut above that inference, once for each premise & every new rank-\(r\) Cut has smaller height; duplicated internal Cuts have rank below \(r\) \\
monotonicity introduces the cut occurrence & discard the vacuous branch and enlarge the other premise to the conclusion & removes the selected summand \\
\(B\land C\), \(B\lor C\), or \(B\to C\) is principal on both sides & use the displayed rule premises to make one or two Cuts on \(B\) and \(C\) & every replacement Cut has rank below \(r\); duplicate contexts collapse as sets \\
an initial \(\bot L\) is encountered & close the retained-\(\bot\) conclusion immediately, or commute on the opposite premise when the Cut formula is that \(\bot\) & removes the selected summand or reduces its height \\
\bottomrule
\end{tabularx}
\end{center}

\begin{hbtheorem}{thm:p07-prop-cut-termination}{proved=annex:p7.arithmetic.1}{Termination of the finite-set propositional Cut strategy}
One step of the strategy in
Definition~\ref{def:p07-prop-cut-measure} preserves the endsequent and strictly
decreases \(o(D)\).  Consequently the strategy terminates, and every finite
derivation in the finite-set support presentation can be transformed into a
cut-free derivation of the same endsequent.  By the support translations, every
propositional derivation on \texttt{sheet:lk-defneg} likewise has a cut-free
derivation of the same multiset endsequent.  The theorem concerns this
deterministic strategy; it does not claim strong normalization of every
possible permutation of Cut reductions, and it does not upgrade the cited
\(\varepsilon_0\)-analysis of arithmetic.
\end{hbtheorem}
% END HISTORICAL SLICE 110-110

% BEGIN HISTORICAL DERIVATIVE 112-309
\begingroup
\renewcommand{\sectionmark}[1]{\markright{\thesection.\ THEORY PRESENTATIONS}}
\section{Theory presentations used in the calibration table}
\endgroup
\label{sec:p07-calibration-theories}

\begin{hbdefinition}{def:p07-pra}{Primitive recursive arithmetic}
The presentation of \emph{primitive recursive arithmetic} used here, written
\(\mathsf{PRA}\), has a function symbol for every primitive-recursive
numerical function, the corresponding defining equations, classical
equational logic, and quantifier-free induction.  Equivalently presented
free-variable systems may replace this surface only after their equivalence
has been stated.  This definition names a formal first-order theory; it is not
a definition of a philosophical standpoint.  Its ordinal calibration is used
in the first row of Table~\ref{tab:p07-ordinal-calibration}.
\index{primitive recursive arithmetic}
\index{PRA|see{primitive recursive arithmetic}}
\end{hbdefinition}

\begin{hbdefinition}{def:p07-heyting-arithmetic}{Heyting arithmetic}
\emph{Heyting arithmetic}, \(\mathsf{HA}\), has the arithmetic language,
nonlogical axioms, equality axioms, and the Peano-arithmetic induction scheme
fixed in Section~\ref{sec:arithmetic-coding-01}, but uses intuitionistic
first-order logic instead of classical first-order logic.  It is a formal
first-order theory, not a semantic notion of constructive truth.  This
distinction is used by the calibration row here and by every soundness and
derivability statement in Chapter~\ref{sec:p07-functional-interpretations}.
\index{Heyting arithmetic}
\index{HA|see{Heyting arithmetic}}
\end{hbdefinition}

\begin{hbdefinition}{def:p07-isigma-fragments}{The fragments \(\mathsf I\Sigma_n\)}
For each \(n\geq1\), the first-order theory \(\mathsf I\Sigma_n\) is
Robinson arithmetic as fixed in Section~\ref{sec:arithmetic-coding-01} together with
exactly the induction instances whose induction formula lies in the
arithmetical class \(\Sigma^0_n\).  This is a family of formal theories indexed
by \(n\), not one unindexed theory and not the syntactic class
\(\Sigma^0_n\) itself.  The distinction is used in the fragment row of
Table~\ref{tab:p07-ordinal-calibration}.
\index{induction fragment!ISigma-n@\(\mathsf I\Sigma_n\)}
\end{hbdefinition}

\begin{hbdefinition}{def:p07-second-order-calibration-theories}{The second-order calibration theories}
Work in the two-sorted language of second-order arithmetic.  The base theory
\(\mathsf{RCA}_0\) uses classical two-sorted logic, the number-sort axioms of
Robinson arithmetic as fixed in Section~\ref{sec:arithmetic-coding-01}, extensionality
for set variables, induction for \(\Sigma^0_1\) formulas, and comprehension
for \(\Delta^0_1\) formulas.  The theory
\(\mathsf{ACA}_0\) is \(\mathsf{RCA}_0\) plus comprehension for every
arithmetical formula.  The theory \(\mathsf{ATR}_0\) is
\(\mathsf{RCA}_0\) plus the arithmetical-transfinite-recursion scheme along
each coded well-order: every arithmetically defined stage operator has a set
coding its stagewise iteration along that order.

These are axiomatically introduced second-order theories, not alternative
names for \(\PA\) or \(\mathsf{HA}\).  Their set-existence schemes distinguish
them from the first-order theories above, and the distinction is used only in
the two corresponding rows of Table~\ref{tab:p07-ordinal-calibration}
\cite{Simpson2009}.
\index{reverse mathematics!RCA0@\(\mathsf{RCA}_0\)}
\index{reverse mathematics!ACA0@\(\mathsf{ACA}_0\)}
\index{reverse mathematics!ATR0@\(\mathsf{ATR}_0\)}
\end{hbdefinition}

\section{Relative ordinal calibrations}
\label{sec:p07-ordinal-calibration}

\begin{hbdefinition}{def:p07-relative-proof-theoretic-ordinal}{A formula-class-relative proof-theoretic ordinal}
Fix a coding that effectively enumerates recursive binary relations
\(e\mapsto(D_e,\prec_e)\), where each \(D_e\) is a recursive subset of
\(\N\), and fix a class \(\mathcal C\) of formulas.  For a theory \(T\), one
standard first-order calibration is the supremum of
\(\operatorname{otp}(D_e,\prec_e)\) over those \(e\in\N\) for which
\((D_e,\prec_e)\) is externally a well-order and \(T\) proves every instance
of \(\mathcal C\)-transfinite induction along \(\prec_e\).
The index varies across the set whose supremum is taken; fixing one relation
and then reusing it in the set-builder would not define this calibration.  An
order type enters the set only after external well-foundedness of its coded
relation has been established.  This definition names a metalinguistic
calibration relative to the formula class, relation coding, and induction
reading; it is not an intrinsic ordinal object inside \(T\).  The
qualifications are nonredundant: consistency, reflection, provably total
functions, and provable well-foundedness can yield different comparisons.
The table below uses this distinction rather than an unqualified vertical-bar
notation for a theory; vertical bars are reserved for cardinality.
\index{proof-theoretic ordinal!relative calibration}
\index{transfinite induction!ordinal analysis}
\end{hbdefinition}

\begingroup\small
\begin{longtable}{@{}p{.22\textwidth}p{.19\textwidth}p{.50\textwidth}@{}}
\caption{Local citation-qualified ordinal-analysis calibrations for the six
displayed theory presentations.  This is input to, not a substitute for, the
separate master table.}
\label{tab:p07-ordinal-calibration}\\
\toprule
Theory/presentation & Calibrated ordinal & Qualification and source \\
\midrule
\endfirsthead
\toprule
Theory/presentation & Calibrated ordinal & Qualification and source \\
\midrule
\endhead
Modern \(\mathsf{PRA}\), with all primitive-recursive function symbols and
quantifier-free induction & \(\omega^\omega\) & Conventional first-order
induction/provably-total-function calibration; it does not define finitism
\cite{Pohlers2009} \\
\(\mathsf I\Sigma_n\), fixed \(n\geq1\) & \(\omega_{n+1}\), where
\(\omega_0=1\) and \(\omega_{k+1}=\omega^{\omega_k}\) & Standard
fragment calibration for the displayed indexing convention
\cite{Pohlers2009} \\
\(\PA\) & \(\varepsilon_0=\sup_n\omega_n\) & Gentzen
consistency/induction calibration, relative to the notation and metatheory in
Section~\ref{sec:normalization-curry-howard-04}
\cite{Gentzen1936,TroelstraSchwichtenberg2000} \\
Heyting arithmetic \(\mathsf{HA}\) & \(\varepsilon_0\) & Standard ordinal
analysis comparison; equality of the displayed calibration does not identify
the classical and intuitionistic proof systems \cite{Pohlers2009} \\
\(\mathsf{ACA}_0\) & \(\varepsilon_0\) & Second-order calibration under the
standard recursive-well-order convention; not a claim that
\(\mathsf{ACA}_0=\PA\) \cite{Simpson2009,Pohlers2009} \\
\(\mathsf{ATR}_0\) and the named autonomous predicative progressions &
\(\Gamma_0\) & The reverse-mathematical and autonomous-progression
calibrations must remain distinguished even when they display the same
ordinal \cite{Feferman1964,Schuette1977,Simpson2009} \\
\bottomrule
\end{longtable}
\endgroup

The table deliberately omits weaker first-order arithmetics, set theories,
and large-cardinal extensions: no number is guessed for them.  The separate
master apparatus must fix and source its lower endpoint before placing
Robinson arithmetic or bounded-induction theories on the same scale; the local
table makes no claim to discharge that broader range.

\section{Veblen notation and the Feferman--Schuette boundary}
\label{sec:p07-veblen}

\begin{hbdefinition}{def:p07-veblen-hierarchy}{Veblen hierarchy and \(\Gamma_0\)}
An ordinal function is \emph{normal} when it is strictly increasing and is
continuous at every nonzero limit.  Set
\(\operatorname{Veb}_0(\beta)=\omega^\beta\).  For \(\alpha>0\), let
\(\operatorname{Veb}_\alpha\) enumerate in increasing order the common fixed
points of all \(\operatorname{Veb}_\gamma\) with \(\gamma<\alpha\).  The
ordinal \(\Gamma_0\) is
the least positive \(\gamma\) satisfying
\(\operatorname{Veb}_\gamma(0)=\gamma\).

The locally defined symbol \(\operatorname{Veb}\) names this external
hierarchy of ordinal functions; it is distinct from the global formula
metavariables and is needed in the displayed definition of \(\Gamma_0\).
These clauses also define the one ordinal \(\Gamma_0\).  The hierarchy extends
the Cantor-normal-form vocabulary of Part~6 because
fixed points of earlier normal functions must themselves be enumerated.  The
distinction is used in the final row of
Table~\ref{tab:p07-ordinal-calibration}; no ordinal-analysis claim beyond
\(\Gamma_0\) is made here.
\index{Veblen hierarchy}
\index{Gamma-zero@\(\Gamma_0\)}
\end{hbdefinition}

\begin{hbtheorem}{thm:p07-feferman-schuette}{cited=Feferman1964}{Feferman--Schuette predicative calibration}
In the autonomous-progression analyses of Feferman and Schuette,
\(\Gamma_0\) is the bound reached by that formalization of predicative
analysis.
\end{hbtheorem}

The statement is about the cited formal autonomous progressions
\cite{Feferman1964,Schuette1977}; it is not the unqualified philosophical
identity ``predicative mathematics is exactly everything below
\(\Gamma_0\).''  Veblen's original hierarchy is \cite{Veblen1908}; a modern
ordinal-analysis reference is \cite{Pohlers2009}.

\section{Primitive recursive arithmetic and finitism}
\label{sec:p07-pra-finitism}

Skolem's 1923 recursive arithmetic is a historical precursor
\cite{Skolem1923}.  Tait argues that finitistically acceptable numerical
functions and arguments are captured by primitive-recursive methods
\cite{Tait1981}.  The latter is Tait's analysis of finitism, not a theorem
that all philosophical accounts identify finitism with the formal theory
\(\mathsf{PRA}\).

\begin{hbwarning}{def:p07-pra-finitism-boundary}{PRA is not a definition of finitism}
``Finitism'' names a methodological and philosophical class of reasoning;
\(\mathsf{PRA}\) is the exact formal theory in
Definition~\ref{def:p07-pra}.  Skolem's historical program and Tait's later
analysis motivate a comparison, but neither licenses replacing the two names
by an identity sign.  In particular, the unrelated bibliography key
\texttt{Skolem1934} is not evidence for the 1923 arithmetic source.
\index{finitism}
\end{hbwarning}

\begin{hbdefinition}{def:p07-omega-calculus}{The selected semiformal arithmetic calculus}
Let \(\mathcal L_2\) have number and set variables, membership, \(0,S\), and a
symbol for every primitive-recursive numerical function and relation.  An
\emph{arithmetic formula} has no set quantifier; it is \emph{semi-closed} when
it has no free number variable.  Positive and negative literals are evaluated
using the standard numerical meanings of their closed terms.  Let
\(\mathsf{PA}_{\infty}\) be exactly the one-sided Tait calculus on
\hyperref[sheet:omega-arithmetic]{\texttt{sheet:omega-arithmetic}}.  Its finite
sequent sets are read disjunctively, and its derivations are well-founded,
possibly countably branching trees.

The name is necessary because neither finitary \(\mathsf{PA}\) nor an arbitrary
``omega logic'' determines these sequents and initials.  It names a proof
system, differs from finitary \(\mathsf{PA}\) by its universal-number rule and
semiformal trees, and is used by the height-and-cut-rank judgment below.  A
finitary induction instance embeds by deriving each numeral instance from the
base and step cases and applying the displayed omega rule
\cite{Schuette1977,JagerProbst2004}.
\index{omega rule@\(\omega\)-rule}\index{omega logic@\(\omega\)-logic}
\end{hbdefinition}

\begin{hbdefinition}{def:p07-omega-derivability}{Height and cut rank in \(\mathsf{PA}_{\infty}\)}
For a literal \(E\), put \(\operatorname{rn}(E)=0\); put
\(\operatorname{rn}(A\circ B)=\max(\operatorname{rn}(A),
\operatorname{rn}(B))+1\) for \(\circ\in\{\land,\lor\}\), and
\(\operatorname{rn}(QxA)=\operatorname{rn}(A)+1\) for a number quantifier
\(Q\).  For an ordinal \(\alpha\), a natural number \(n\), and a finite
sequent \(\Gamma\), define
\(\mathsf{PA}_{\infty}\vdash_n^\alpha\Gamma\) by induction on \(\alpha\):

\begin{enumerate}
  \item if \(\Gamma\) is one of the two displayed initials, then the judgment
        holds for every \(\alpha,n\);
  \item if every premise \(\Gamma_i\) of one displayed rule satisfies
        \(\mathsf{PA}_{\infty}\vdash_n^{\alpha_i}\Gamma_i\) for some
        \(\alpha_i<\alpha\), then its conclusion satisfies the judgment,
        provided that every cut formula has rank strictly below \(n\).
\end{enumerate}

Thus the superscript bounds tree height and the subscript bounds cut-formula
rank; subscript \(0\) means cut-free.  Finally define
\(2_0(\alpha)=\alpha\) and \(2_{m+1}(\alpha)=2^{2_m(\alpha)}\).  This relation is
needed because an unindexed assertion of infinitary derivability cannot state
the quantitative cut-elimination theorem.  It is a metalinguistic relation on
derivations, not an arithmetic predicate; its two indices respectively differ
from formula rank and from the ordinal value alone, and both occur in the
theorem below \cite[Definition~1 and Theorem~2]{JagerProbst2004}.
\index{cut rank!omega arithmetic@\(\omega\)-arithmetic}
\end{hbdefinition}

\begin{hbtheorem}{thm:p07-omega-cut-elimination}{cited=JagerProbst2004}{Cut elimination and the \(\varepsilon_0\) boundary for the selected \(\omega\)-calculus}
For every finite sequent \(\Gamma\), ordinal \(\alpha\), and
\(n\in\mathbb N\),
\[
  \mathsf{PA}_{\infty}\vdash_n^\alpha\Gamma
  \quad\Longrightarrow\quad
  \mathsf{PA}_{\infty}\vdash_0^{2_n(\alpha)}\Gamma .
\]
This is J\"ager and Probst's Theorem~2 for the exact calculus and indexed
derivability relation displayed above.  A finite \(\mathsf{PA}\) proof embeds
at some finite cut rank and some height below \(\varepsilon_0\); because
\(\varepsilon_0\) is closed under \(\alpha\mapsto2^\alpha\), the resulting
cut-free height remains below \(\varepsilon_0\).  The statement imports the
proved reduction theorem with its exact domain and bound; it does not pretend
that a bare ordinal or an unspecified local reduction table proves termination
\cite{Schuette1977,Pohlers2009,JagerProbst2004}.
\end{hbtheorem}

The chapter therefore keeps four distinct claims separate: termination for the
selected finite propositional reduction, the Veblen/Feferman--Schuette
calibration, the formal PRA versus philosophical-finitism boundary, and the
cited cut-elimination theorem for the exact \(\omega\)-rule calculus.
% END HISTORICAL DERIVATIVE 112-309
